% =============================================================
%  main.tex — KDD AAE 2026 Workshop Submission
%  "A Governed Harness for Auditable LLM-Driven ML Experimentation"
%
%  Overleaf setup:
%    1. Upload this file and references.bib to the same Overleaf project.
%    2. Set the main document to main.tex.
%    3. Compiler: pdfLaTeX (default).
%    4. Bibliography style: natbib + plainnat (included below).
%
%  KDD 2026 uses the ACM SIG proceedings template.
%  Replace \documentclass below with the official KDD template class
%  once the workshop releases it. The structure below is compatible.
% =============================================================

\documentclass[sigconf]{acmart}

% -- Packages --------------------------------------------------
\usepackage{booktabs}       % \toprule, \midrule, \bottomrule
\usepackage{microtype}      % better justification
\usepackage{hyperref}       % clickable links
\usepackage{url}            % \url{}
\usepackage{listings}       % pseudocode blocks
\usepackage{graphicx}       % \includegraphics
\usepackage{multirow}       % table row spanning

% -- ACM metadata (update before submission) -------------------
\acmConference[KDD AAE'26]{KDD Workshop on Agentic AI Evaluation}{August 2026}{Barcelona, Spain}
\acmDOI{}
\acmISBN{}
\copyrightyear{2026}

% =============================================================
\begin{document}

\title{A Governed Harness for Auditable LLM-Driven ML Experimentation}

\author{[Author names omitted for review]}
\affiliation{\institution{[Institution omitted for review]}}
\email{[omitted]}

% =============================================================
\begin{abstract}
Autonomous ML experimentation agents can propose code edits, execute
training runs, and select follow-up trials, but their evaluation remains
difficult because conventional task metrics do not capture invalid actions,
repeated failures, auditability, or lifecycle governance.
We present a governed harness for autonomous ML experimentation.
In the Agent~=~Model~+~Harness decomposition~\cite{trivedy2026anatomy},
our harness owns trial lifecycle, scope enforcement, and audit; the LLM
manager is a replaceable component.
The framework separates proposal, execution, validation, and decision
authority; enforces editable-scope constraints; records append-only trial
ledgers; and exports governance metrics including invalid-action rate,
repeated-bad rate, artifact completeness, and provenance completeness.
We demonstrate the framework on a ResNet-trigger scientific ML benchmark
with a real full-loop campaign that executes bounded code edits and records
complete provenance.
We further define a memory-ablation protocol---comparing no memory, summary
memory, and rationale-linked append-only memory under fixed budget---grounded
in harness engineering practice~\cite{trivedy2026betterharness,bockeler2026coders}.
The study argues that agentic experimentation should be evaluated not only by
final task performance, but also by whether agent behavior is bounded,
auditable, failure-aware, and reproducible.
\end{abstract}

\keywords{agentic AI, autonomous experimentation, harness engineering,
          governed control plane, audit, reproducibility, memory ablation}

\maketitle

% =============================================================
\section{Introduction}
\label{sec:introduction}

% TODO: write from plan/KDD_AAE_refinement_plan_v2.md §9.4 hook.
% Lead: autonomous agents can plan experiments, edit code, execute runs —
% but evaluating by final score alone is insufficient.
% Introduce the control plane as the answer.

[Introduction text here.]

% =============================================================
\section{Background and Related Work}
\label{sec:related}

% Source: paper/kdd_aae_2026/sections/02_related_work.md

\subsection{Autonomous Code-Space Exploration}

AIDE~\cite{jiang2025aide} is the closest related academic system.
Both AIDE and \texttt{autoresearch\_harness} use LLMs to propose atomic
code changes evaluated by a hard-coded metric.
The structural difference is where lifecycle authority lives: AIDE couples
search policy, evaluator, and summarisation inside one optimisation loop,
while this work separates proposal generation from governance.

% Table 5 (AIDE comparison) goes here — see plan §7.3 for columns.

SHARP~\cite{birk2026sharp} frames scientific reproduction as a translation
task and demonstrates human-agent collaboration for reproducing a
particle-physics result.
SHARP targets faithful reproduction with human oversight at checkpoints;
this work targets governed autonomous experimentation without per-trial
human intervention.
These are complementary contributions.

The AI Scientist~\cite{lu2024aiscientist} targets open-ended autonomous
discovery.
Mind2Web~\cite{deng2023mind2web} constructs broad web-agent benchmarks
across 137 domains.
Unlike benchmark-driven evaluation across many domains, we prioritise a
single real task node with complete provenance, trading breadth for
inspectability.

\subsection{Harness Engineering}

Designing environments, feedback loops, and control systems is now
recognised as the primary engineering challenge for agentic
AI~\cite{bockeler2026firstthoughts,openai2026codex}.
Trivedy~\cite{trivedy2026anatomy} defines the harness as everything around
the model: tools, sandboxes, memory, middleware, and execution hooks.
B\"{o}ckeler~\cite{bockeler2026coders} argues that trust requires constraining
the solution space with guides and sensors.
OpenAI~\cite{openai2026codex} describes production agent work as an
observability and feedback-loop problem: give the agent a map, not a manual.
Anthropic~\cite{anthropic2025longrunninga,anthropic2026longrunnningb}
emphasises progress tracking, verifier separation, and skeptical evaluation
for long-running agents.
Hashimoto~\cite{hashimoto2026journey} independently derived the same
principle: anytime an agent makes a mistake, engineer a structural fix so
it never recurs.

\subsection{Experiment Tracking and AutoML}

Experiment trackers such as MLflow~\cite{mlflow2018} and Weights \&
Biases~\cite{wandb2020} record outcomes and metrics after a run exists.
They do not govern the agent loop, enforce bounded execution before a patch
is applied, or classify invalid trials as first-class audit objects.
AutoML systems optimise within a predefined search space but are not
designed to audit self-directed agent exploration or explain why a proposed
code change was kept, discarded, or failed invalid.

% =============================================================
\section{Governed Control Plane}
\label{sec:system}

% Source: paper/kdd_aae_2026/sections/03_system_design.md
% Key subsections: Manager/Worker Separation, Trial Lifecycle,
%   Scope Enforcement, Pending-Trial Guard, Append-Only Ledger,
%   Memory Modes, Decision Authority, Guides/Sensors table,
%   Failure Mode Motivation table.

[System design text here.]

% Guides/Sensors 2x2 table (Table S1):
\begin{table}[h]
\caption{Harness guides and sensors for autonomous ML experimentation,
  instantiating the taxonomy of~\citet{bockeler2026coders}.}
\label{tab:guides-sensors}
\small
\begin{tabular}{p{1.6cm}p{2.8cm}p{2.8cm}}
\toprule
 & \textbf{Guide (feedforward)} & \textbf{Sensor (feedback)} \\
\midrule
\textbf{Computational} &
  Node spec; editable-path whitelist; budget cap; scope constraints &
  Metric parser; state machine; pending-trial guard \\
\midrule
\textbf{Inferential} &
  Manager system prompt; memory injection; proposal rationale template &
  Repeated-bad detector (edit-target and mechanism similarity) \\
\bottomrule
\end{tabular}
\end{table}

% Failure Mode Motivation table (Table S2):
\begin{table}[h]
\caption{Agent failure modes mapped to ML experimentation analogues and
  harness controls, adapted from~\citet{anthropic2025longrunninga}.}
\label{tab:failure-modes}
\small
\begin{tabular}{p{2.0cm}p{2.2cm}p{2.2cm}}
\toprule
\textbf{Failure mode} & \textbf{ML analogue} & \textbf{Harness fix} \\
\midrule
Declares victory early  & Stops after one good trial      & Budget enforcer \\
Leaves broken state     & Failed patch dirties repo       & Pending-trial guard \\
Marks done prematurely  & Claims improvement without metric & \texttt{FAILED\_INVALID} + metric parser \\
Cannot run the app      & Out-of-scope edit proposed      & Node spec + editable-path whitelist \\
\bottomrule
\end{tabular}
\end{table}

% =============================================================
\section{Evaluation Protocol}
\label{sec:experiments}

% Source: paper/kdd_aae_2026/sections/04_experiments.md
% Key subsections: Benchmark Node, Fixed-Budget Campaign Protocol,
%   Memory Ablation Design (pre-stated hypothesis), Stress Trial,
%   Governance Metrics, Failure Taxonomy (Table 2).

% Memory ablation hypothesis (must appear BEFORE results):
We hypothesise, consistent with~\citet{openai2026codex}
and~\citet{bockeler2026coders}, that structured failure context
(rationale memory) reduces the repeated-bad rate more than raw
trial summaries, because it gives the manager a targeted failure
signal rather than history noise:
\[
  \text{RepeatedBad}(\texttt{none}) >
  \text{RepeatedBad}(\texttt{summary}) >
  \text{RepeatedBad}(\texttt{summary+rationale})
\]

[Experiments text here.]

% =============================================================
\section{Results}
\label{sec:results}

% Source: paper/kdd_aae_2026/sections/05_results.md
% Order: governance lifecycle demo -> memory ablation ->
%   failure taxonomy -> provenance completeness -> task metric.

[Results text here.]

% =============================================================
\section{Discussion and Limitations}
\label{sec:discussion}

% Source: paper/kdd_aae_2026/sections/06_discussion_limitations.md

\paragraph{Single benchmark node.}
All experiments use the ResNet-trigger node.
We cannot rule out that reported improvements overfit to this evaluation
domain.
Applying the harness hill-climbing methodology of~\citet{trivedy2026betterharness}
across multiple evaluation nodes with holdout splits would strengthen the
governance claims; this is future work.

\paragraph{Self-evaluation authority.}
We keep the keep/discard decision in the control plane rather than
delegating to the manager.
\citet{anthropic2026longrunnningb} found that LLM self-evaluation is
systematically lenient; a deterministic, externally-owned decision criterion
addresses this.

\paragraph{Non-claims.}
This work does not claim to build a general autonomous scientist, prove
scientific discovery, introduce a universal optimisation algorithm, or
depend on a specific coding-agent backend.
The ResNet-trigger task is a real scientific ML case study for evaluating
governed autonomous experimentation; it is not claimed to represent all ML
optimisation tasks.

% =============================================================
\section{Conclusion}
\label{sec:conclusion}

[Conclusion text here.]

% =============================================================
\bibliographystyle{ACM-Reference-Format}
\bibliography{references}

\end{document}
